\definecolor{hl1}{RGB}{51,51,51}
\definecolor{hl2}{RGB}{24,54,145}
\definecolor{hl3}{RGB}{0,134,179}
\definecolor{hl4}{RGB}{15,0,0}
\definecolor{hl5}{RGB}{3,3,3}
\definecolor{hl6}{RGB}{99,163,92}
\definecolor{hl7}{RGB}{121,93,163}
\definecolor{hl8}{RGB}{167,29,93}
\definecolor{hl9}{RGB}{223,80,0}
\definecolor{hl10}{RGB}{150,152,150}
\emph{Calepin} processes citations from BibTeX/BibLaTeX files using \href{https://github.com/typst/hayagriva}{hayagriva}. Specify one or more \texttt{.bib} files in the front matter, and citation keys in the document body are resolved automatically.

\subsection*{Front matter}
\label{front-matter}

The \texttt{bibliography} field accepts a single path or a list. The optional \texttt{csl} field points to a \href{https://citationstyles.org/}{Citation Style Language} file.




\begin{srccode}
\begin{Verbatim}[commandchars=\\\{\}]
\textcolor{hl1}{---
}
\textcolor{hl1}{title = "My Document"
}
\textcolor{hl1}{bibliography = "references.bib"
}
\textcolor{hl1}{csl = "apa"
}
\textcolor{hl1}{---}
\end{Verbatim}
\end{srccode}

Multiple bibliography files are merged into a single library:




\begin{srccode}
\begin{Verbatim}[commandchars=\\\{\}]
\textcolor{hl1}{---
}
\textcolor{hl1}{bibliography = ["refs\_theory.bib", "refs\_empirical.bib"]
}
\textcolor{hl1}{---}
\end{Verbatim}
\end{srccode}

\subsection*{Citation syntax}
\label{citation-syntax}

Three citation forms are supported:

\medskip
\begin{tblr}{colspec={Q[l]Q[l]X[l]}, hline{1,2,Z}={solid}}
Syntax & Output & Description \\
\texttt{@key} & Author et al. (Year) & Narrative citation \\
\texttt{{[}@key{]}} & (Author et al. Year) & Parenthetical citation \\
\texttt{{[}-@key{]}} & Year & Suppress author, year only \\
\texttt{{[}@key1; @key2{]}} & (Author1 Year1; Author2 Year2) & Grouped citations \\
\end{tblr}
\medskip

Grouped references also work for cross-references: \texttt{{[}\hyperref[fig-one]{Figure 1}; \hyperref[fig-two]{Figure 2}{]}} produces linked references like {[}\hyperref[fig-one]{Figure 1}; \hyperref[fig-two]{Figure 2}{]}.

These are compatible with Quarto and Pandoc syntax.

\subsection*{CSL styles}
\label{csl-styles}

By default, \emph{Calepin} uses the Chicago author-date style. To customize citation formatting, set the \texttt{csl} field in the front matter or in \texttt{\{stem\}\_calepin/config.toml}. The value can be a built-in style name or a file path:




\begin{srccode}
\begin{Verbatim}[commandchars=\\\{\}]
\textcolor{hl1}{---
}
\textcolor{hl1}{csl = "apa"                          \# built-in name (no .csl extension)
}
\textcolor{hl1}{csl = "styles/my-journal.csl"       \# file path
}
\textcolor{hl1}{---}
\end{Verbatim}
\end{srccode}

Resolution order (first match wins):

\begin{enumerate}
\item File path --- If the \texttt{csl} value points to an existing \texttt{.csl} file, it is used directly.


\item Built-in name --- \emph{Calepin} ships with hundreds of styles from the \href{https://github.com/typst/hayagriva}{hayagriva} archive. Run \texttt{calepin extra csl} to list all available names.


\item Default --- If no \texttt{csl} field is set, \emph{Calepin} falls back to \texttt{chicago-author-date}.


\end{enumerate}

If the specified style is not found or is a dependent style, \emph{Calepin} warns and falls back to the default.

Additional \texttt{.csl} files for journals not covered by the built-in archive are available from the \href{https://www.zotero.org/styles}{Zotero style repository}.

\subsection*{References section}
\label{references-section}

When a document contains at least one resolved citation, \emph{Calepin} appends a References section at the end of the output automatically. Unresolved keys (those not found in any \texttt{.bib} file) are left as-is.

Two dummy figures so we can demonstrate cross-references:



\begin{figure}
\centering
\fbox{\parbox[c][50.0pt]{50.0pt}{\centering 100×100}}


\caption{First figure}
\label{fig-one}
\end{figure}

\begin{figure}
\centering
\fbox{\parbox[c][50.0pt]{50.0pt}{\centering 100×100}}


\caption{Second figure}
\label{fig-two}
\end{figure}